\documentclass[A4,10.5pt]{article}
% \documentclass[letterpaper,11pt]{article} %For use in US
\usepackage{latexsym}
\usepackage[empty]{fullpage}
\usepackage{titlesec}
\usepackage{marvosym}
\usepackage[usenames,dvipsnames]{color}
\usepackage{verbatim}
\usepackage{enumitem}
\usepackage[colorlinks=true, linkcolor=blue, urlcolor=NavyBlue, citecolor=blue]{hyperref}
\usepackage[english]{babel}
\usepackage{tabularx}
\usepackage{tikz}
\input{glyphtounicode}

\usepackage{palatino}
\usepackage{tabularx} 
\addtolength{\oddsidemargin}{-0.6in}
\addtolength{\evensidemargin}{-0.6in}
\addtolength{\textwidth}{1.2in}
\addtolength{\topmargin}{-0.8in}
\addtolength{\textheight}{1.7in}

\urlstyle{same}

\definecolor{datecolor}{HTML}{2E6DA4}

\raggedbottom
\raggedright
\setlength{\tabcolsep}{0cm}

\titleformat{\section}{  \vspace{-4pt}\scshape\raggedright\large\bfseries
}{}{0em}{}[\color{black}\titlerule \vspace{-4pt}]


\pdfgentounicode=1
\newcommand{\CVItem}[1]{
  \item\small{
    {#1 \vspace{-2pt}}
  }
}

\newcommand{\CVSubheading}[4]{
  \vspace{3pt}\item
    \begin{tabular*}{0.97\textwidth}[t]{p{0.8\textwidth}@{\extracolsep{\fill}}r}
      \textbf{#1} & \small\textcolor{datecolor}{#2} \\
      \textit{\small #3} & \textit{\small\textcolor{datecolor}{#4}} \\
    \end{tabular*}\vspace{-4pt}
}

\newcommand{\CVSubSubheading}[2]{
    \item
    \begin{tabular*}{0.97\textwidth}{l@{\extracolsep{\fill}}r}
      \text{\small#1} & \text{\small #2} \\
    \end{tabular*}\vspace{-7pt}
}

\newcommand{\CVProject}[2]{
  \vspace{3pt}\item
    \begin{tabular*}{0.97\textwidth}[t]{p{0.9\textwidth}@{\extracolsep{\fill}}r}
      \textbf{#1} & #2 \\
      % \small#3 & \\
    \end{tabular*}\vspace{0pt}
}

\newcommand{\CVWorkExp}[2]{
  \vspace{-2pt}\item
    \begin{tabular*}{0.97\textwidth}[t]{p{0.8\textwidth}@{\extracolsep{\fill}}r}
      \textbf{#1} & #2 \\
      % \small#3 & \\
    \end{tabular*}
}



\newcommand{\CVEduItem}[4]{
  \vspace{-2pt}\item
    \begin{tabular*}{0.97\textwidth}[t]{l@{\extracolsep{\fill}}r}
      \textbf{#1} $|$ \small\textit{#3}, #4 & \small\textcolor{datecolor}{#2} \\
    \end{tabular*}\vspace{-15pt}
}

\newcommand{\CVSubItem}[1]{\CVItem{#1}\vspace{-4pt}}

\renewcommand\labelitemii{$\vcenter{\hbox{\tiny$\bullet$}}$}

\newcommand{\CVSubHeadingListStart}{\begin{itemize}[leftmargin=0.5cm, label={}]}
% \newcommand{\resumeSubHeadingListStart}{\begin{itemize}[leftmargin=0.15in, label={}]} % Uncomment for US
\newcommand{\CVSubHeadingListEnd}{\end{itemize}}
\newcommand{\CVItemListStart}{\begin{itemize}}
\newcommand{\CVItemListEnd}{\end{itemize}\vspace{-11pt}}

%------------------------------------------------------------------------------
% CV STARTS HERE  %
%------------------------------------------------------------------------------
\begin{document}


%==================== HEADER ====================%
\begin{center}
    \textbf{\huge Vivek Sai} \\ \vspace{2pt}
    \small
    +49 155 10421564 \;|\;
    \href{mailto:viveksaisurya07@gmail.com}{\underline{viveksaisurya07@gmail.com}} \;|\;
    \href{https://www.linkedin.com/in/vivek-sai-5b363221a/}{\underline{LinkedIn}} \;|\;
    \href{https://github.com/VivekSai07}{\underline{GitHub}} \;|\;
    \href{https://vivek-portfolio-mocha.vercel.app/}{\underline{Portfolio}}
\end{center}

\vspace{-8pt}

\section{Education}
  \CVSubHeadingListStart
    \CVEduItem{M.Sc. Computer Science (Autonomous Systems)}{Oct 2024 -- Present}{University of Stuttgart}{Stuttgart, Germany}
    \CVEduItem{B.Tech. Computer Science \& Engineering (AI)}{Oct 2020 -- Apr 2024}{Amrita Vishwa Vidyapeetham}{Chennai, India}
  \CVSubHeadingListEnd

%-----WORK EXPERIENCE----------------------------------------------------------
\vspace{4pt}
\section{Work Experience}
\CVSubHeadingListStart
\CVSubheading
  {Working Student -- Humanoid Robotics, Development \& Validation}{Apr 2026 -- Present}
  {Porsche Engineering Services GmbH}{Stuttgart, Germany}
\CVItemListStart
    \CVItem{Deployed and validated an NVIDIA cuVSLAM + nvblox perception-localization stack on Unitree Go2/G1 (Jetson Orin NX, RealSense D435i), resolving DDS, TF-tree, and domain-ID conflicts to sustain 30~Hz visual odometry and real-time 3D reconstruction.}
    \CVItem{Unified Go2 and G1 onto a single arm64 YOLOv8/TensorRT detection pipeline (Isaac ROS), fixing build, QoS, and runtime defects to reach real-time detections at $\sim$24~Hz on embedded Jetson hardware.}
    \CVItem{Designed a SOLID-refactored ROS~2 service converting 2D object detections into 3D poses at 30~Hz with full unit-test coverage, enabling downstream navigation and manipulation.}
    \CVItem{Re-architected the robot audio stack: unified text-to-speech into a shared cross-robot scheduler and shipped a swappable ASR backend (Nemotron) with VAD-based turn detection and wake-word gating.}
\CVItemListEnd
\CVSubheading
  {Student Assistant}{Nov 2025 -- Jun 2026}
  {Socially Intelligent Robotics (SIR) Lab, University of Stuttgart}{Stuttgart, Germany}
\CVItemListStart
    \CVItem{Debugged hardware--software integration issues on 2 FER robotic platforms, cutting experiment setup failures by $\sim$40\% and improving uptime.}
    \CVItem{Prototyped and evaluated learning-based robotic behaviors in NVIDIA Isaac Sim, focusing on sim-to-real transfer.}
\CVItemListEnd
\CVSubheading
  {Student Assistant}{Dec 2024 -- Mar 2026}
  {Construction Robotics, University of Stuttgart}{Stuttgart, Germany}
\CVItemListStart
    \CVItem{Implemented and validated pose estimation and control loops, conducting sim-to-real validation that identified a $\sim$15\% perception accuracy gap and reduced it to under 5\% through calibration and pipeline tuning.}
    \CVItem{Trained a PPO-based reinforcement learning policy for peg-in-hole insertion on a Franka arm in simulation, tuning reward shaping, action space, and hyperparameters to improve policy convergence.}
    \CVItem{Trained a SmolVLA (VLA) model on a SO-100 arm in NVIDIA Isaac Sim, achieving 40--60\% inference success.}
\CVItemListEnd
\CVSubheading
  {Research Intern}{Jun 2023 -- Aug 2023}
  {IIITDM Kancheepuram}{Chennai, India}
\CVItemListStart
    \CVItem{Implemented and benchmarked deep learning models (CNNs, Vision Transformers) in PyTorch/TensorFlow, with data preprocessing, augmentation, and evaluation pipelines to analyze robustness and generalization.}
\CVItemListEnd
\CVSubHeadingListEnd

%-----PROJECTS AND RESEARCH----------------------------------------------------
% \vspace{0.5pt}
\section{Projects}
\CVSubHeadingListStart
    \CVProject
      {\href{https://github.com/VivekSai07/ContactPilot}{ContactPilot} — Vision-Based Robotic Grasping Pipeline $|$ \emph{\small{Python, PyTorch, MuJoCo, Open3D, SAM 3, LLMs}}}{}
    \CVItemListStart
        \CVItem{Built an end-to-end pick-and-place pipeline for a Franka Panda arm with two interchangeable 6-DoF grasp backends (Contact-GraspNet, NVIDIA GraspGen), multi-camera point-cloud fusion, and vision-only occupancy-based bin placement for collision-free packing.}
        \CVItem{Designed a seed-based benchmarking harness with automated failure-taxonomy analysis, raising pick-and-place success from 38\% to a sustained 100\% (30/30 objects, 10-seed benchmark); added SAM 3 promptable object selection (mask IoU 0.29 $\to$ 0.89--0.99) and an LLM instruction parser (Llama 3.1 via NVIDIA NIM) for multi-step free-text commands.}
    \CVItemListEnd
    \CVProject
      {\href{https://github.com/VivekSai07/franka-vla-generalization}{Generalizable VLA Policies for Franka Manipulation} $|$ \emph{\small{Python, MuJoCo, robosuite, HuggingFace}}}{}
    \CVItemListStart
        \CVItem{Designed a systematic multi-axis generalization study (spatial, language, distractors) fine-tuning SmolVLA and X-VLA on sealed, scripted datasets (24--560 episodes); published 7 datasets and 9 model checkpoints to Hugging Face.}
        \CVItem{X-VLA reached 88.1\% on unseen colour-instruction pairs (held-out within 10 pts of trained) and 94.4\% on language grounding with no positional shortcut verified; identified and fixed a classifier bug via video inspection.}
    \CVItemListEnd
    \CVProject
      {\href{https://github.com/VivekSai07/amr_odom}{3D LiDAR Odometry} (MAD-ICP) $|$ \emph{\small{C++, ROS~2, Eigen, Python}}}{}
    \CVItemListStart
        \CVItem{Implemented a MAD-ICP-inspired 3D LiDAR odometry system in C++ (Eigen-only core + ROS~2 Humble wrapper); source-voxel downsampling (0.15~m) gave $\sim$1.9$\times$ wall-clock speedup on 128-beam LiDAR, achieving ATE RMSE 0.014~m (handheld), RPE drift 2.8\%/4.4\% on SPOT loop/stair datasets.}
    \CVItemListEnd
\CVSubHeadingListEnd
    

\vspace{-6pt}
\section*{Skills}
\begin{tabular}{p{0.20\textwidth} p{0.80\textwidth}}
\textbf{Programming:} & C++, Python, Linux/Bash \\
\textbf{Robotics:} & ROS~2, robot kinematics, motion execution, path planning, localization, sim-to-real transfer \\
\textbf{Perception:} & Intel RealSense D435i, NVIDIA cuVSLAM, nvblox, camera calibration, 6D pose estimation (AprilTag) \\
\textbf{Simulation:} & NVIDIA Isaac Sim, MuJoCo, Gazebo, URDF \\
\textbf{ML / Deep Learning:} & PyTorch, CNNs, Vision Transformers, reinforcement learning (PPO), VLA fine-tuning (SmolVLA, X-VLA), model evaluation \\
\textbf{Hardware:} & Unitree Go2 (quadruped), Unitree G1 (humanoid), Jetson Orin NX, FER robotic platforms, SO-100 \\
\textbf{Tools:} & Git, Docker, unitree\_sdk2 \\
\textbf{Languages:} & English (C1), German (A1)
\end{tabular}




% \vspace{2pt}
% \section{Journal Publications}
% \CVSubHeadingListStart
%   \item \small \href{https://link.springer.com/article/10.1007/s42979-025-03862-0}{HybridViT: An Approach for Alzheimer’s Disease Classification with ADNI Neuroimaging Data} \hfill Apr 2025 \\
%   \textit{Springer Nature Computer Science}
% \CVSubHeadingListEnd


\end{document}